% Fonts, layout, links
\usepackage[margin=1in]{geometry}
\usepackage{setspace}
\onehalfspacing
\usepackage{lmodern}
\usepackage[T1]{fontenc}
\usepackage[utf8]{inputenc}
\usepackage[spanish]{babel}
\usepackage[hidelinks]{hyperref}
% \hypersetup{
%   pdftitle={\ThesisTitle},
%   pdfauthor={\AuthorName}
% }

% Figures, tables, math
\usepackage{graphicx}
\graphicspath{{Fig/}}
\usepackage{caption}
\usepackage{subcaption}
\usepackage{booktabs}
\usepackage{array}
\usepackage{longtable}
\usepackage{tabularx}
\usepackage{amsmath,amssymb,amsfonts}
\usepackage{siunitx}
\usepackage{float}
\usepackage{xcolor}
\usepackage{listings}
\usepackage{svg}

% listing formatting
\lstset{
    basicstyle=\ttfamily\small,
    frame=single,
    numbers=left,
    numberstyle=\tiny\color{gray},
    numbersep=10pt,
    breaklines=true,
    breakatwhitespace=false,
    showstringspaces=false,
    tabsize=4,
    columns=fullflexible,
    keepspaces=true,
    xleftmargin=2em,
    framexleftmargin=2em
}

% Colores
\definecolor{fondo}{RGB}{248,248,248}
\definecolor{marco}{RGB}{210,210,210}
\definecolor{palabrasclave}{RGB}{120,60,150}
\definecolor{tipos}{RGB}{30,120,90}
\definecolor{numeros}{RGB}{140,140,140}
\definecolor{funciones}{RGB}{180,100,30}
\definecolor{comentarios}{RGB}{100,100,100}

% Lenguaje para el modelo relacional
\lstdefinelanguage{relacional}{
    morekeywords={PRIMARY,KEY,FOREIGN,REFERENCES},
    morekeywords=[2]{CHAR,VARCHAR,INT,DATE,TIME,FLOAT,DECIMAL,BOOLEAN},
    sensitive=false
}

% Lenguaje SQL
\lstdefinelanguage{sql}{
    morekeywords={
        SELECT,FROM,WHERE,INSERT,INTO,VALUES,UPDATE,SET,DELETE,
        CREATE,ALTER,DROP,TABLE,DATABASE,INDEX,VIEW,
        JOIN,INNER,LEFT,RIGHT,FULL,OUTER,ON,
        GROUP,BY,ORDER,ASC,DESC,HAVING,LIMIT,OFFSET,
        DISTINCT,AS,AND,OR,NOT,IN,IS,NULL,BETWEEN,LIKE,
        UNION,ALL,EXISTS,
        PRIMARY,KEY,FOREIGN,REFERENCES,CONSTRAINT,
        UNIQUE,CHECK,DEFAULT,AUTO_INCREMENT,
        CASCADE,RESTRICT
    },
    morekeywords=[2]{
        CHAR,VARCHAR,TINYINT,SMALLINT,INT,BIGINT,
        FLOAT,DOUBLE,DECIMAL,DATE,DATETIME,TIME,
        TIMESTAMP,YEAR,BOOLEAN,TEXT,TINYTEXT,MEDIUMTEXT,LONGTEXT,
        BLOB,TINYBLOB,MEDIUMBLOB,LONGBLOB
    },
    morekeywords=[3]{
        COUNT,SUM,AVG,MIN,MAX,COALESCE,CONCAT,
        UPPER,LOWER,LENGTH,ROUND,NOW,CURDATE,CURTIME
    },
    sensitive=false,
    morecomment=[l]{--},
    morecomment=[s]{/*}{*/},
    morestring=[b]'
}

% Estilo para el modelo relacional
\lstdefinestyle{modelo}{
    language=relacional,
    basicstyle=\ttfamily\small,
    backgroundcolor=\color{fondo},
    frame=single,
    rulecolor=\color{marco},
    framesep=8pt,
    numbers=left,
    numberstyle=\tiny\color{numeros},
    numbersep=10pt,
    keywordstyle=\color{palabrasclave}\bfseries,
    keywordstyle=[2]\color{tipos},
    breaklines=true,
    keepspaces=true,
    showstringspaces=false,
    tabsize=4
}

% Estilo para SQL
\lstdefinestyle{codigoSQL}{
    language=sql,
    basicstyle=\ttfamily\small,
    backgroundcolor=\color{fondo},
    frame=single,
    rulecolor=\color{marco},
    framesep=8pt,
    numbers=left,
    numberstyle=\tiny\color{numeros},
    numbersep=10pt,
    keywordstyle=\color{palabrasclave}\bfseries,
    keywordstyle=[2]\color{tipos},
    keywordstyle=[3]\color{funciones}\bfseries,
    commentstyle=\color{comentarios}\itshape,
    stringstyle=\color{tipos},
    breaklines=true,
    keepspaces=true,
    showstringspaces=false,
    tabsize=4
}




% Chapter/section formatting
\usepackage{titlesec}
\titleformat{\chapter}[block]
  {\normalfont\bfseries\Large\centering} % formato del título
  {}{0pt}{\MakeUppercase} % etiqueta vacía: no muestra "CAPÍTULO N:"
\titlespacing*{\chapter}{0pt}{-0.5em}{1.5em}
\setcounter{secnumdepth}{3}
\setcounter{tocdepth}{3}

% Number figures/tables by chapter
\usepackage{chngcntr}
\counterwithin{figure}{chapter}
\counterwithin{table}{chapter}

% Include LOF/LOT/References in ToC
\usepackage[nottoc]{tocbibind}
\usepackage{tikz}
\usetikzlibrary{arrows.meta,positioning,shapes,calc,fit}
\usepackage{csquotes}
% Citation setup (edit style if needed)
\usepackage[backend=biber,style=ieee,sorting=none,maxnames=6]{biblatex}
\addbibresource{bib/references.bib}

% Handy macros for consistent names
\newcommand{\UnivLine}{\textbf{\University}}
